% amsmath documentation, chapter 'Displayed equations'
\documentclass[oneside]{book}

\title{BookML test: \texttt{amsmath} user guide, displayed equations}
\author{American Mathematical Society, \LaTeX\ Project}
\date{1999-12-13\\(revised 2002-02-25, 2016-11-14, 2018-04-05, 2019-10-14, 2020-02-18)}

\usepackage{amsmath}

\usepackage{hyperref}

\hypersetup{pdfborder=0 0 0}

\usepackage{bookml/bookml}
\usepackage{framed}
\AtBeginDocument{\begin{framed}
  \begin{description}
    \item[BookML] \BookMLversion
    \item[\LaTeXML] \LaTeXMLversion{} (\csname bml@generator@engine\endcsname)
    \item[\LaTeX] \fmtversion
  \end{description}
\end{framed}}


\begin{document}

\chapter{Displayed equations}

\section{Introduction}

\begin{equation*}
a=b
\end{equation*}

\begin{equation}
a=b
\end{equation}

\begin{equation}\label{xx}
\begin{split}
a& =b+c-d\\
 & \quad +e-f\\
 & =g+h\\
 & =i
\end{split}
\end{equation}

\begin{multline}
a+b+c+d+e+f\\
+i+j+k+l+m+n
\end{multline}

\begin{gather}
a_1=b_1+c_1\\
a_2=b_2+c_2-d_2+e_2
\end{gather}

\begin{align}
a_1& =b_1+c_1\\
a_2& =b_2+c_2-d_2+e_2
\end{align}

\begin{align}
a_{11}& =b_{11}&
  a_{12}& =b_{12}\\
a_{21}& =b_{21}&
  a_{22}& =b_{22}+c_{22}
\end{align}

\begin{flalign*}
a_{11} + b_{11}& = c_{11}&
  a_{12}& =b_{12}\\
b_{21}& = c_{21}&
  a_{22}& =b_{22}+c_{22}
\end{flalign*}

\section{Split equations without alignment}

\begin{multline}
\framebox[.65\columnwidth]{A}\\
\framebox[.5\columnwidth]{B}\\
\shoveright{\framebox[.55\columnwidth]{C}}\\
\framebox[.65\columnwidth]{D}
\end{multline}

\section{Split equations with alignment}

\begin{equation}\label{e:barwq}\begin{split}
H_c&=\frac{1}{2n} \sum^n_{l=0}(-1)^{l}(n-{l})^{p-2}
  \sum_{l _1+\dots+ l _p=l}\prod^p_{i=1} \binom{n_i}{l _i}\\
&\quad\cdot[(n-l )-(n_i-l _i)]^{n_i-l _i}\cdot
  \Bigl[(n-l )^2-\sum^p_{j=1}(n_i-l _i)^2\Bigr].
  \kern-2em % adjust equation body to the right [mjd,13-Nov-1994]
\end{split}\end{equation}

\section{Equation groups without alignment}

\begin{gather}
  first equation\\
  \begin{split}
    second & equation\\
           & on two lines
  \end{split}
  \\
  third equation
\end{gather}

\section{Equation groups with mutual alignment}

\begin{align}
x&=y       & X&=Y       & a&=b+c\\
x'&=y'     & X'&=Y'     & a'&=b\\
x+x'&=y+y' & X+X'&=Y+Y' & a'b&=c'b
\end{align}

\begin{align}
x& = y_1-y_2+y_3-y_5+y_8-\dots
                    && \text{by \eqref{eq:C}}\\
 & = y'\circ y^*    && \text{by \eqref{eq:D}}\\
 & = y(0) y'        && \text {by Axiom 1.}
\end{align}

\begin{alignat}{2}
x& = y_1-y_2+y_3-y_5+y_8-\dots
                  &\quad& \text{by \eqref{eq:C}}\\
 & = y'\circ y^*  && \text{by \eqref{eq:D}}\\
 & = y(0) y'      && \text {by Axiom 1.}
\end{alignat}

\begin{flalign}
x&=y       & X&=Y\\
x'&=y'     & X'&=Y'\\
x+x'&=y+y' & X+X'&=Y+Y'
\end{flalign}

\begin{flalign*}
x&=y       & X&=Y\\
x'&=y'     & X'&=Y'\\
x+x'&=y+y' & X+X'&=Y+Y'
\end{flalign*}

\section{Alignment building blocks}

\begin{equation*}
\left.\begin{aligned}
  B'&=-\partial\times E,\\
  E'&=\partial\times B - 4\pi j,
\end{aligned}
\right\}
\qquad \text{Maxwell's equations}
\end{equation*}

\begin{equation}\label{eq:C}
P_{r-j}=
  \begin{cases}
    0&  \text{if $r-j$ is odd},\\
    r!\,(-1)^{(r-j)/2}&  \text{if $r-j$ is even}.
  \end{cases}
\end{equation}

\section{Interrupting a display}

\begin{align}
  A_1&=N_0(\lambda;\Omega')-\phi(\lambda;\Omega'),\\
  A_2&=\phi(\lambda;\Omega')-\phi(\lambda;\Omega),\\
\intertext{and}
  A_3&=\mathcal{N}(\lambda;\omega).
\end{align}


\chapter{Miscellaneous mathematical features}

\section{Matrices}

\begin{equation}\label{eq:D}
\begin{pmatrix} D_1t&-a_{12}t_2&\dots&-a_{1n}t_n\\
-a_{21}t_1&D_2t&\dots&-a_{2n}t_n\\
\hdotsfor[2]{4}\\
-a_{n1}t_1&-a_{n2}t_2&\dots&D_nt\end{pmatrix},
\end{equation}

\chapter{Additional BookML tests}

\section{Subequations}

\begin{subequations}
  \begin{align}
    a_1& =b_1+c_1\\
    a_2& =b_2+c_2-d_2+e_2\\
    a_{11}& =b_{11}&
      a_{12}& =b_{12}\\
    a_{21}& =b_{21}&
      a_{22}& =b_{22}+c_{22}
  \end{align}
\end{subequations}

\end{document}
